% !TeX spellcheck = pt_BR
\documentclass[a4paper,12pt]{letter}  % tipo do documento
\usepackage[brazil]{babel}           % define a linguagem padrão
\usepackage{graphicx}                % permite a inserção de figuras
\usepackage[T1]{fontenc}
\usepackage[utf8]{inputenc}
\usepackage[hmargin=1.5cm,vmargin=1cm]{geometry}
\usepackage{listings} %para colocar codigos
\usepackage{amsmath}
\usepackage{amssymb}
\usepackage{multicol}
\setlength{\columnsep}{1.5cm}
\setlength{\columnseprule}{0.2pt}

% \newcommand{\gabarito}[1]{\textbf{\begin{small} \\ #1 \\ \end{small}}} %mostrar gabarito
\newcommand{\gabarito}[1]{} %nao mostrar gabarito


\begin{document}
\thispagestyle{empty}

UNIVERSIDADE FEDERAL DO CEARÁ\\
Departamento de Estatística e Matemática Aplicada\\
Disciplina de Elementos de Análise Combinatória CC0279\\
Professor: Albert Einstein F. Muritiba\\
Temas: Análise Combinatória\\[0.2cm]
\textbf{Nome}: \makebox[380pt]{\hrulefill} \vspace{2pt}
%Matr\'icula:\makebox[60pt]{\hrulefill} \vspace{2pt}
\begin{center}
	2ª Avaliação Parcial (AP)
	-- \textit{2ª Chamada}
	% 2\textordfeminine Avaliação - Segunda Chamada
	% npc 2 %+ nef
	% NEF
\end{center}

\begin{enumerate}
	% \item (2pts.) Defina:
	%       \begin{enumerate}
	% 	      \item Permutação (simples, com repetição, circular, caótica);
	% 	      \item Combinação (simples, com repetição);
	%       \end{enumerate}
	\item Para que servem os lemas de Kaplansky? Qual a diferença entre os dois lemas?
	\item Uma rede de \textit{fast-food} disponibiliza 16 ingredientes sólidos (tomate, cebola, etc.), 4 tipos de molho e 4 tipos de macarrão. Calcule de quantas formas é possível montar um prato com 8 ingredientes sólidos, 2 molhos e 1 tipo de macarrão, sabendo que a escolha de ingredientes e molhos pode ser repetida.
	% \item Qual o número mínimo de pessoas em um grupo para que possamos garantir que nele haja pelo menos 10 pessoas nascidas no mesmo mês?
	% \item De quantos modos é possível enfileirar \textbf{h} homens e \textbf{m} mulheres, todos de alturas distintas, de modo que os homens, entre si, e as mulheres, entre si, estejam em ordem crescente de altura?
	% \item De quantos modos podemos escolher três dias para estudar cálculo semanalmente, sem escolher dias consecutivos?
	% \item (2pts.)De quantas maneiras é possível colocar 6 anéis em 4 dedos (da mão) já predefinidos, sendo que:
	%       \begin{enumerate}
	% 	      \item os anéis são \textbf{todos diferentes} e a ordem dos anéis nos dedos \textbf{importa};
	% 	      \item os anéis são \textbf{todos diferentes}, mas a ordem dos anéis nos dedos \textbf{não importa};
	% 	      \item os anéis são \textbf{iguais};
	% 	      \item os anéis são \textbf{iguais}, porém os dedos \textbf{não estão predefinidos}.
	%       \end{enumerate}
	\item De quantas maneiras é possível colocar 6 anéis em 4 dedos já predefinidos, sendo que os anéis são \textbf{iguais}?

	\item Quantos números no intervalo [1, 1000] são divisíveis por pelo menos um dos números 2, 3 e 5?
	% \item Quantos números no intervalo [1, 1000] são divisíveis por pelo menos dois dos números 2, 3 e 5?

	\item Considere o conjunto $A = \{1, 2, 3, 4, 5, 6, 7, 8, 9\}$. De quantas formas podemos ordenar seus elementos de modo que exatamente 3 elementos estejam em suas posições originais?


\end{enumerate}


\begin{scriptsize}
	\textbf{Fórmulas:}\\
	$P_n = n!$\\
	$P^{a,b, ...}_n = \frac{n!}{a!b! ...}$\\
	$PC_n = (n-1)!$\\
	$C^p_n = \frac{n!}{p!(n-p)!}$\\
	$CR^{p}_{n} = C^{p}_{n+p-1}$\\
	Kaplansky: $f(n,p) = C^p_{n-p+1}$ , $g(n,p) = \frac{n}{n-p}C^p_{n-p}$ \\
	Permutação caótica: $D_n = n! \left[ \frac{1}{0!} - \frac{1}{1!} + \frac{1}{2!} - \frac{1}{3!} + \dots + (-1)^n \frac{1}{n!}\right] $\\
	Princípio da Inclusão-Exclusão:\\
	$\sharp \bigcup_{i=1}^n A_i = \sum_{i=1}^n (-1)^{i-1} S_i$\\
	$S_0 = \sharp\Omega $ ,
	$S_1 = \sum_{i=1}^n \sharp A_i $ ,
	$S_2 = \sum_{1\leq i < j \leq n}^n \sharp(A_i \cap A_j) $ ,
	$S_3 = \sum_{1\leq i < j < k \leq n}^n \sharp(A_i \cap A_j \cap A_k) $ \dots \\
	Exatamente p: $\sum_{k=0}^{n-p} (-1)^k C^{k}_{k+p} S_{p+k}$\\
	Pelo menos p: $\sum_{k=0}^{n-p} (-1)^k C^{k}_{k+p-1} S_{p+k}$\\
\end{scriptsize}
%\end{multicols}
\end{document}
